% Gerado por scripts/gerar_integracao.py.
\begin{frame}{Integração: de milhares a bilhões de transistores}
\only<1|handout:0>{\figura{\graficoIntegracao{1}}\etapa{1/2}{Marcos Intel: mais transistores por circuito integrado.}}
\only<2|handout:1>{\figura{\graficoIntegracao{2}}\etapa{2/2}{SoCs Apple e GPU NVIDIA ampliam o panorama; funções e áreas são distintas.}}
\fonte{Fontes primárias: Intel (guia histórico e Nehalem), ISSCC 2012 (Ivy Bridge), Apple Newsroom (M1--M4) e NVIDIA (GA100). Eixo vertical logarítmico. Contagens por die, não por placa; processo é a designação publicada, não comprimento de porta. Dados e links em dados/integracao\_processadores.csv.}
\end{frame}
\begin{frame}{Ler o gráfico: o que significa ``integrar mais?}
\duas{\textcolor{blue_theme}{\large\textbf{O que contamos}}\par\medskip Mais dispositivos no mesmo circuito integrado.\par\medskip CPU, memórias, interfaces e aceleradores podem dividir o die.}{\textcolor{blue_theme}{\large\textbf{O que também importa}}\par\medskip Processo, área do die e função do chip.\par\medskip A contagem sozinha não mede desempenho nem densidade.}
\chamada{VLSI: integração de muitos dispositivos e conexões em um chip.}
\fonte{M1--M4 são SoCs com arquitetura Arm; A100 é uma GPU. O histórico é uma seleção didática de produtos, sem ajuste de uma curva universal de Moore. As contagens recentes abrangem o chip completo.}
\end{frame}
